% Generated draft: requires mathematical and source review.
\documentclass[11pt]{article}
\usepackage[margin=25mm]{geometry}
\usepackage{fontspec}
\setmainfont{Latin Modern Roman}
\usepackage{amsmath,amssymb,mathtools,unicode-math}
\setmathfont{Latin Modern Math}
\usepackage{enumitem,needspace,xurl,hyperref}
\setlength{\parindent}{0pt}
\setlength{\parskip}{5pt}
\setlength{\emergencystretch}{3em}
\newcommand{\N}{\mathbb{N}}
\newcommand{\Z}{\mathbb{Z}}
\newcommand{\Q}{\mathbb{Q}}
\newcommand{\R}{\mathbb{R}}
\newcommand{\C}{\mathbb{C}}
\newcommand{\F}{\mathbb{F}}
\newcommand{\A}{\mathbb{A}}
\newcommand{\PP}{\mathbb{P}}
\newcommand{\E}{\mathbb{E}}
\newcommand{\eps}{\varepsilon}
\newcommand{\dd}{\,\mathrm{d}}
\begin{document}
\section*{UnsolvedMath ID: 700810}
\textbf{Source problem code:} catalogue entry
\subsection{Definitions and mathematical statement}
All groups considered are finite. A formation is a nonempty class \(\mathfrak F\) of groups, closed under isomorphisms and homomorphic images, with the following subdirect-product closure property: if \(N_1,\ldots,N_r\trianglelefteq G\), \(r\geq 1\), \(\bigcap_{i=1}^r N_i=1\), and \(G/N_i\in\mathfrak F\) for every \(i\), then \(G\in\mathfrak F\). It is hereditary if every subgroup of each member also belongs to \(\mathfrak F\). It is soluble if every member is soluble, where solubility means that the derived series terminates at the identity subgroup.

A chief factor of \(G\) is a quotient \(A/B\), where \(B<A\) are normal subgroups of \(G\) and no normal subgroup of \(G\) lies strictly between them. Write
\[
C_G(A/B)=\{g\in G: gag^{-1}a^{-1}\in B\text{ for every }a\in A\}.
\]
Using the standard chief-factor convention, a formation \(\mathfrak F\) is local if there is a function \(f\), assigning to each prime \(p\) a formation or the empty class, such that, for every finite group \(G\),
\[
G\in\mathfrak F
\quad\Longleftrightarrow\quad
G/C_G(A/B)\in f(p)
\quad\text{for every chief factor }A/B\text{ of }G
\text{ and every prime }p\mid |A/B|.
\]
For the trivial group the condition on chief factors is vacuous.

An \(\mathfrak F\)-covering subgroup of \(G\) means a subgroup \(U\leq G\) such that \(U\in\mathfrak F\) and
\[
\forall H\ (U\leq H\leq G)\quad
\forall N\trianglelefteq H,\qquad
H/N\in\mathfrak F\ \Longrightarrow\ H=UN.
\]
Here \(UN=\{un:u\in U,\ n\in N\}\).

\textbf{Classification problem (L. A. Shemetkov, 1990).} Determine all hereditary soluble local formations \(\mathfrak F\) satisfying
\[
\forall\text{ finite groups }G\quad
\exists U\leq G\quad
\left[
U\in\mathfrak F\ \land\
\forall H\ (U\leq H\leq G)\ 
\forall N\trianglelefteq H\,
\bigl(H/N\in\mathfrak F\Longrightarrow H=UN\bigr)
\right].
\]
The quantifier over \(G\) includes all finite groups, including nonsoluble groups; solubility is a requirement on the members of \(\mathfrak F\). The supplied record does not explicitly define the technical terms; the standard conventions specified above require confirmation against the cited source during mathematical review.
\subsection{Short English statement}
Classify the hereditary soluble local formations \(\mathfrak F\) for which every finite group has an \(\mathfrak F\)-covering subgroup. Such a subgroup belongs to \(\mathfrak F\) and maps onto every quotient belonging to \(\mathfrak F\) of every intermediate subgroup containing it.
\subsection{Sources}
\begin{itemize}
\item Local UnsolvedMath registry, permanent ID \(700810\), exact source problem code \texttt{catalogue entry}, title \emph{Kourovka Notebook Problem 11.118}. The supplied registry status is \texttt{open}; this is a supplied label, not an independently verified assessment.
\item \emph{The Kourovka Notebook (Unsolved Problems in Group Theory)}, original Problem 11.118, 11th issue (1990); attributed in the supplied record to L. A. Shemetkov, with proposed year 1990.
\item The supplied background identifies the 21st issue edition (2026), \texttt{arXiv:1401.0300v48}. The supplied collection description names E. I. Khukhro and V. D. Mazurov as editors of the Issues 1--21 collection (1965--2026). These bibliographic details are reproduced from the record without independent source verification.
\end{itemize}
\end{document}
